% Part: proof-theory
% Chapter: proof-search
% Section: tableaux

\documentclass[../../../include/open-logic-section]{subfiles}

\begin{document}

\olfileid{pt}{ps}{tab}

\olsection{في الجداول الدلالية}

أنظمة الجداول الدلالية من أنظمة~!!{proof} القريبة من حسابات المتتاليات، وهي
صالحة بخاصة للبحث عن البرهان يدويًا أو بالبرمجيات. غير أن~!!a{proof} فيها
ليس شجرة متتاليات، بل شجرة من~!!{formula}s؛ وقواعدها، خلافًا للاستنباط
الطبيعي، صورة معكوسة لقواعد حساب المتتاليات. و!!{proof} الجدول الدلالي في
حقيقته بحث صريح عن نموذج، ومن ثم تسمى هذه الجداول أيضًا \emph{أشجار صدق}.

\emph{الجدول الدلالي الموقع} شجرة من~!!{formula}s ذات علامات، كل منها على
الصورة~$\sFmla{\True}{!A}$ أو~$\sFmla{\False}{!A}$. وتنمو الشجرة إلى
أسفل، ومبدؤها~!!{formula}s تعين المطلوب أن~!!{prove}. فإذا أردنا، مثلًا،
أن~!!{prove} المتتالية~$!A,!B\Sequent !C,!D$، بدأنا بـ
$\sFmla{\True}{!A},\sFmla{\True}{!B},
\sFmla{\False}{!C},\sFmla{\False}{!D}$. وكل~!!^a{structure} يجعل~$!A$
و$!B$ صادقتين، و$!C$ و$!D$ كاذبتين، فهو~!!a{structure} يشهد بأن
$!A,!B\Sequent !C,!D$ غير صالحة.

وتوسع العقدة الورقية، أي السفلى، في جدول دلالي بقواعد تمديد الفروع. فتزيد
هذه القواعد عقدًا في الفرع نفسه، أو تزيد تحته عقدتين شقيقتين فتشعبه فرعين.
ويجوز إعمالها في كل~!!{formula} موقعة فوق العقدة في الفرع نفسه. وقواعد
تمديد الفروع لنظام الجداول الدلالية الكلاسيكي~\Log{Tc} مثبتة في
\olref{tab:Tc}.

\subfile{rules-Tc}

وما هذه القواعد، كما ترى، إلا قواعد~\Log{G3c} مقلوبة؛ فالعلامتان~$\True$
و$\False$ تدلان، على الترتيب، على أن~!!a{formula} رئيسة عن يسار المتتالية
أو عن يمينها.

ونسمي فرع جدول دلالي \emph{مغلقًا} إذا اشتمل معًا على
$\sFmla{\True}{!A}$ و$\sFmla{\False}{!A}$، أو احتوى
$\sFmla{\True}{\lfalse}$. فإن أغلقت الفروع كلها كان الجدول الدلالي نفسه
مغلقًا. ومن أمثلته جدول دلالي مغلق للمتتالية
$!A\lor(!B\land !C)\Sequent(!A\lor !B)\land(!A\lor !C)$؛ أي إنه يبدأ بـ
$\sFmla{\True}{!A\lor(!B\land !C)}$ و
$\sFmla{\False}{(!A\lor !B)\land(!A\lor !C)}$:
\begin{maxwidth}
\begin{tableau}{}
  [\sFmla{\True}{\formula{A} \lor (\formula{B} \land \formula{C})},
    just=\TAss, checked
    [\sFmla{\False}{(\formula{A} \lor \formula{B}) \land
        (\formula{A} \lor \formula{C})}, just=\TAss, checked
      [\sFmla{\True}{\formula{A}}, just={\TRule{\True}{\lor}[1]}
        [\sFmla{\False}{\formula{A} \lor \formula{B}},
          just={\TRule{\False}{\land}[2]}, checked
          [\sFmla{\False}{\formula{A}},
            just={\TRule{\False}{\lor}[4]}
            [\sFmla{\False}{\formula{B}},
              just={\TRule{\False}{\lor}[4]}, close]
          ]
        ]
        [\sFmla{\False}{\formula{A} \lor \formula{C}},
          just={\TRule{\False}{\land}[2]}, checked
          [\sFmla{\False}{\formula{A}},
            just={\TRule{\False}{\lor}[4]}
            [\sFmla{\False}{\formula{C}},
              just={\TRule{\False}{\lor}[4]}, close]
          ]
        ]
      ]
      [\sFmla{\True}{\formula{B} \land \formula{C}},
        just={\TRule{\True}{\lor}[1]}, checked
        [\sFmla{\False}{\formula{A} \lor \formula{B}},
          just={\TRule{\False}{\land}[2]}, checked
          [\sFmla{\True}{\formula{B}},
            just={\TRule{\True}{\land}[3]}, move by=2
            [\sFmla{\True}{\formula{C}},
              just={\TRule{\True}{\land}[3]}
              [\sFmla{\False}{\formula{A}},
                just={\TRule{\False}{\lor}[4]}
                [\sFmla{\False}{\formula{B}},
                  just={\TRule{\False}{\lor}[4]},close
                ]
              ]
            ]
          ]
        ]
        [\sFmla{\False}{\formula{A} \lor \formula{C}},
          just={\TRule{\False}{\land}[2]}, checked
          [\sFmla{\True}{\formula{B}},
            just={\TRule{\True}{\land}[3]}, move by=2
              [\sFmla{\True}{\formula{C}},
                just={\TRule{\True}{\land}[3]}
                [\sFmla{\False}{\formula{A}},
                  just={\TRule{\False}{\lor}[4]}
                  [\sFmla{\False}{\formula{C}},
                  just={\TRule{\False}{\lor}[4]},close
                ]
              ]
            ]
          ]
        ]
      ]
    ]
  ]
\end{tableau}
\end{maxwidth}

وعلامة الصح تدل على إعمال القاعدة المقابلة في~!!a{formula} بجميع الفروع
التي تحتها، وأما~$\otimes$ فتدل على إغلاق الفرع. وأرقام الأسطر بعد
القواعد تعين~!!{formula}s الموقعة التي عملت فيها القاعدة، أي~!!{formula}
الواقعة في الفرع نفسه في السطر المشار إليه.

والبحث عن البرهان في الجداول الدلالية يسير. فننشئ الجدول الدلالي على مراحل: في المرحلة
${\OLMathNumeral{0}}$ نضع~!!{formula}s الابتدائية في أعلاه، وفي المرحلة~${\OLId{latin-ordinary}{n}}+{\OLMathNumeral{1}}$ نأخذ كل عقدة
ورقية في فرع مفتوح، فنجد أول صيغة موقعة غير ذرية، من جهة الأعلى، لم تؤشر،
ونعمل فيها قاعدة تمديد الفرع. فإذا عملت القاعدة في كل فرع مفتوح مشتمل على
تلك الصيغة أشرناها. وهذا المنع من التكرار مقصور على الصيغ غير الذرية في
الفروع المفتوحة. (وبهذه الخوارزمية ولد المثال السابق.)

وأما~!!{formula}s الموقعة من الصورتين
$\sFmla{\True}{\lforall}$ و$\sFmla{\False}{\lexists}$ فلا تؤشر قط،
ولذلك تفرد بمعالجة خاصة. فإذا وجدت وجب أن نضمن تطبيق
$\TRule{\True}{\lforall}$ و$\TRule{\False}{\lexists}$ في النهاية على جميع
الحدود المنطبقة. ويمكن ضمان ذلك، مثلًا، بتطبيقهما في كل مرحلة على جميع هذه
الصيغ في كل فرع مفتوح، بعد معالجة أعلى~!!{formula} موقعة لا تكون من هاتين
الصورتين. والحد~${\OLId{latin-ordinary}{t}}$ المستعمل هو أول حد يمكن بناؤه من~!!{constant}s الواقعة
بالفعل في الفرع (أو~$\Obj c_{\OLMathNumeral{0}}$ إذا لم يقع أي~!!{constant}) ومن
!!{function}s الواقعة في~!!{formula}s الابتدائية، ولا تقع الصيغة الموقعة
المقابلة~!!{formula}~$\sFmla{\True}{!B({\OLId{latin-ordinary}{t}})}$ أو~$\sFmla{\False}{!B({\OLId{latin-ordinary}{t}})}$ في الفرع من
قبل. وإذا لم يوجد حد منطبق غير مستعمل، كانت الصيغة المسوّرة مستنفدة في تلك
المرحلة، فلا نضيف لها عقدة؛ ونعود إليها في مرحلة لاحقة إذا أدخل الفرع
ثوابت جديدة فظهرت حدود منطبقة جديدة.

فإن لم ينتج البحث جدولًا دلاليًا مغلقًا، فإما أن يكون فيه فرع مفتوح منتهٍ
قد أُشّرت فيه جميع~!!{formula}s غير الذرية سوى الصيغ من الصورتين اللتين
لا تؤشّران، واستُنفدت لكل واحدة من هذه الصيغ المسوّرة جميع الحدود القابلة
للتكوين من~!!{constant}s الفرع و!!{function}s الواقعة في~!!{formula}s
الابتدائية؛ وإما أن ينشئ شجرة لامتناهية متناهية التشعب، فيلزم أن يكون
لها فرع لامتناهٍ. وكما في
\olref[cpl]{sec}، نعرف~!!a{structure}~$\Struct{M}$ بحيث متى وقعت
$\sFmla{\True}{!A}$ في الفرع كان~$\Sat{{\OLId{latin-looped}{M}}}{!A}$، وإذا وقعت
$\sFmla{\False}{!A}$ كان~$\Sat/{{\OLId{latin-looped}{M}}}{!A}$. (خذ
$\Theta=\Setabs{!A}{\sFmla{\True}{!A}\text{ في الفرع}}$ و
$\Xi=\Setabs{!A}{\sFmla{\False}{!A}\text{ في الفرع}}$.)

وتنقل الجداول الدلالية المنشأة بهذا الوجه بغير عسر إلى براهين في~\Log{G3c}،
بأن تسند إلى كل عقدة متتالية. فإذا قابلت الصيغ الابتدائية المتتالية
${\OLId{greek-variable}{Gamma}}\Sequent{\OLId{greek-variable}{Delta}}$، فوسم كل صيغة ابتدائية بـ
${\OLId{greek-variable}{Gamma}}\Sequent{\OLId{greek-variable}{Delta}}$. وإذا مُدِّد فرع بتطبيق قاعدة تمديد على~!!{formula}
موقعة~$\sFmla{\True}{!A}$، فوسم الورقة بـ
$!A,\Pi\Sequent\Lambda$؛ وإذا طُبقت على~!!{formula} موقعة
$\sFmla{\False}{!A}$، فوسم الورقة بـ~$\Pi\Sequent\Lambda,!A$. وحيث يعمل
الجدول قاعدة~$\TRule{\True}{\star}$، تعمل قاعدة~$\LeftR{\star}$ في
المتتالية، مع~$!A$ صيغة رئيسة؛ وحيث يعمل~$\TRule{\False}{\star}$، تعمل
قاعدة~$\RightR{\star}$. ومقدمات القاعدة هي المتتاليات الموسومة بها العقد
المقابلة المضافة إلى الجدول. ومتى أغلق جدول بـ
$\sFmla{\True}{!A}$ و$\sFmla{\False}{!A}$، كان وسم العقدة الورقية الأخيرة
متتاليةً من الصورة~$!A,\Pi\Sequent\Lambda,!A$. وإذا أُغلق بـ
$\sFmla{\True}{\lfalse}$، كان وسم الورقة من الصورة
$\lfalse,\Pi\Sequent\Lambda$، وهي مسلّمة~\Log{G3c} كذلك. وقد تُسند
المتتالية نفسها إلى عقد متعاقبة في الجدول؛ فاحذف المكررات. فمثلًا يترجم
الجدول السابق إلى~!!{proof} الآتي:
\begin{maxwidth}
$
\Axiom$!A \fCenter !A, !B$
\RightLabel{\RightR{\lor}}
\UnaryInf$!A \fCenter !A \lor !B$
\Axiom$!A \fCenter !A, !C$
\RightLabel{\RightR{\lor}}
\UnaryInf$!A \fCenter !A \lor !C$
\RightLabel{\RightR{\land}}
\BinaryInf$!A \fCenter (!A \lor !B) \land (!A \lor !C)$
\Axiom$!B, !C \fCenter !A, !B$
\RightLabel{\RightR{\lor}}
\UnaryInf$!B, !C \fCenter !A \lor !B$
\RightLabel{\LeftR{\land}}
\UnaryInf$!B \land !C \fCenter !A \lor !B$
\Axiom$!B, !C \fCenter !A, !C$
\RightLabel{\RightR{\lor}}
\UnaryInf$!B, !C \fCenter !A \lor !C$
\RightLabel{\LeftR{\land}}
\UnaryInf$!B \land !C \fCenter !A \lor !C$
\RightLabel{\RightR{\land}}
\BinaryInf$!B \land !C \fCenter (!A \lor !B) \land
(!A \lor !C)$
\RightLabel{\LeftR{\lor}}
\BinaryInf$!A \lor (!B \land !C) \fCenter (!A \lor !B) \land
(!A \lor !C)$
\DisplayProof$
\end{maxwidth}

\end{document}
